\section{The Natural Exponential -- tmp}
\label{sec:exponential}
\markright{\sc{}The Natural Exponential}

At the most superficial level we use Calculus to create new functions
by defining the new function to be the solution(s) of a specified
\term{Initial Value Problem (IVP)}.

Let's step back for a moment.

We've seen that when we differentiate
the formula $y=x^2$ we get the \term{differential equation} $\d{y} =
2x\d{x}.$ Of course we've abandoned differentials (except as a
``useful fiction'', of course) so we would write this differential
equation as
\begin{equation}
\dfdx{y}{x}=2x.\label{eq:x2diffeq}
\end{equation}
If we did not already know that the function $y(x)=x^2$ satisfies this equation we
could give a name to the solution -- sqr$(x)$, perhaps? --
and by fiat, \alert{define} sqr$(x)$ to be the solution of this
equation.

Oh! But wait! There is a problem. There is more than one solution to
equation~\ref{eq:x2diffeq}, remember?
\begin{ProficiencyDrill}
  Find two more solutions of equation~\ref{eq:x2diffeq}.
\end{ProficiencyDrill}
We can't define sqr$(x)$ to  be
\alert{the} solution if there is more than one solution. This seems
like a big problem until you think about it for a moment. Out of all of
the solutions  of equation~\ref{eq:x2diffeq} only one passes through
the point $(0,0)$ so we will choose to define sqr$(x)$ to be that one.
\begin{ProficiencyDrill}
  Find a formula for sqr$(x)$ that solves each of the following
  IVPs. Graph your solutions:
  \begin{enumerate}[label={\bf{}(\alph*)}]
    \begin{multicols}{2}
    \item $\dfdx{y}{x}=2x,\ \ \ y(0)=0$
    \item $\dfdx{y}{x}=2x,\ \ \ y(0)=1$
    \item $\dfdx{y}{x}=2x,\ \ \ y(0)=10$
    \item $\dfdx{y}{x}=2x,\ \ \ y(1)=2$
    \item $\dfdx{y}{x}=2x,\ \ \ y(-1)=1$
    \end{multicols}
  \end{enumerate}
\end{ProficiencyDrill}

There are many solutions of the differential
equation~\ref{eq:x2diffeq}, but the IVP (Initial Value Problem)
has only one. If we weren't already familiar with the function
$f(x)=x^2$ we could define $\text{sqr}(x)$ to be the solution of
\begin{equation}
\dfdx{y}{x}=2x, \ \ y(0)=0,\label{eq:ExpDiffeq2}
\end{equation}
We can learn a lot about the solution of this IVP just
by examining what the problem itself tells us.

For example since $2x<0$ when $x<0$, the slope of the graph of $y(x)$
is clearly also negative when $x<0$, Similarly, the slope of the graph
of $y(x)$ is positive when $x>0$.

In exactly the same way we can declare, by fiat, the following:\\
\begin{mydefinition}[The Natural Exponential Function]
  \label{def:natural-exponential}
  The function which satisfies the IVP
  \begin{equation}
    \label{eq:ExpDiffeq}
    \dfdx{y}{x}=y, \ \ y(0)=1.
  \end{equation}
  is called the natural exponential  function: $\exp(x)$.
\end{mydefinition}
Traditionally this is called the ``natural''
exponential function.  Of course, simply giving it a name does not
tell us anything about it.

Let's see what we can learn about $\exp(x)$ from
IVP~\ref{eq:ExpDiffeq}.
% Think about the formula $\dfdx{y}{x}=y$ for a moment. It says that if
% $(x,y)$ is a point on the graph of the equation then the slope at that
% point \aside{equivalently, the slope of its tangent line} will be the
% same as the $y$ coordinate.
% % \begin{wrapfigure}[]{l}{1in}
% %   \captionsetup{labelformat=empty}
% %   \centerline{\includegraphics*[height=1in,width=1in]{../Figures/exp-envelope1}}
% %   \label{fig:exp-envelope1}
% % \end{wrapfigure}
% From the initial value we also know that $y(0)=1$ so at the point
% $(0,1)$ the slope of our curve is $1,$ the same as the $y$
% coordinate. Thus near the point $(0,1)$ our curve looks like the line:
% $y=x+1$.

% Think about this for a moment. 
Do you see that the equation
$\dfdx{y}{x}=y$ implies that the curve is always increasing?

\begin{mylemma}
  \label{lemma:ExpInc}
  The $y$ coordinate of every point on the graph of $y=\exp(x)$ is
  always positive and its derivative is also positive. 
\end{mylemma}
\begin{myproof}
  Since $\dfdx{y}{x}=y$ we really only need to show that one of
  $\dfdx{y}{x}$ or $y$ is positive. Think of $x$ as moving along the
  horizontal axis.
  
Since at the point $(0,1)$ the slope is positive, the $y$
  coordinate is rising as the curve passes through it. Thus, once
  the curve has passed $x=0$ the $y$ coordinate must be greater than
  one. Thereafter it continues to increase so $y>1>0$ for every
  positive value of $x.$ Since $\dfdx{y}{x}=y$ the derivative of
  $\exp(x)$ is also positive.

  To show that the $y$ coordinate is positive to the left of $x=0$ as
  well\aside{Careful. We do not mean that moving $x$ to the left
    causes the curve to increase. We mean that the slope of the curve
    is always positive.} is only a little harder.  Supppose that at
  some point, say $x_0$, to the left of $x=0$ the $y$ coordinate is
  negative. Then an argument very similar to the one in the previous
  paragraph would prove that the slope would remain negative to the
  right of $x_0$. But this is impossible since $0$ is to the right of
  $x_0$ and we know that the $y$ coordinate is positive there.

   So $y=\exp(x)$ cannot be negative at $x_0.$ And since we did not
   specify a value for $x_0$ the same argument will work for every
   point to the left of $x=0.$

   Therefore $y=\exp(x)$ is positive and increasing at every $x.$
    \end{myproof}

    We now pause to take explicit notice that something unexpected
    just happened. In Lemma~\ref{lemma:ExpInc} we set out to prove
    that the graph of $y=\exp(x)$ was positive and increasing. We did
    that, but our proof actually tells us a bit more than that. Since
    $y$ is increasing and $\dfdx{y}{x}=y$ we see that the slope of
    $y=\dfdx{y}{x}$ is also increasing. That is, as we move $x$ to the
    right the graph must get steeper and this is always true.
    \begin{ProficiencyDrill}
      What would such a graph look like?
    \end{ProficiencyDrill}
    


It is now time to solve this IVP by finding an explicit  formula for
$\exp(x)$ if we can.

But how? Nothing really presents itself as a potential solution so
what should we do?
Since we don't seem to have any better options, let's see if we can 
guess a solution.
\vskip1mm{}
\begin{Digression}[How to guess effectively]
  No, really. Let's try guessing. It can be surprisingly effective.

  You know this of course. You've been doing it all of your life, but
  in the past -- especially in math class -- you've probably been
  discouraged from guessing so you tend to deny doing it, probably
  even to yourself.

  It's ok. Guess anyway. We won't tell anyone.

  The fact is that guessing is a tried and true solution technique and
  we encourage you to use it regularly. It is especially effective in
  mathematics because mathematics reflects the behavior of the world
  we live in. Just by virtue of being a human being living in the
  world you have a great deal of intuition into that behavior. You
  really do. And you know this. You rely on your intuition almost
  daily to navigate through your life. Guessing is nothing more, or
  less, than relying on your native intuition.

  However, effective guessing is a skill. You need practice to do it
  well and, unfortunately, you have probably not had much
  chance to practice that skill in the context of mathematics.  So we
  encourage you to guess often from now on. But be aware that making a
  guess is a process, not an event. If you
  guess wrong, \underline{\bf{}which is most likely}, then you have a
  new, related problem to think about: All of your intuition said this
  was a good guess.  Why didn't it work?  The answer to that question
  almost always gives some insight into the original problem.  Even a
  bad guess can be useful.

  And bad guesses are easy. If you don't know what else to do, make a
  bad guess. But all by itself a bad guess is a waste of
  time\aside{Indeed, any guess, all by itself, is a waste of
    time.}. You must also take the time to figure out why your bad
  guess doesn't work.

  \underline{\large\bf But} the real danger in guessing -- the reason
  it is usually discouraged -- is that when you do guess correctly, or
  nearly so, it is {\bf very} tempting to just move on from
  there.

  \centerline{\underline{\bf{}Don't do that.}}

  Guessing correctly means that your intuition was very good. But
  intuition is unconcious. If you don't stop and ask yourself why this
  seemed like a good idea your thought process will remain unconscious
  and thus nearly useless. It is at least as important to understand
  where a good guess came from, and why it worked as it is to
  understand why a bad guess didn't work. But when you guess well it
  is extremely tempting to just take your good guess and run with
  it. This will invariably lead to confusion later. So, when you guess
  correctly it is crucial that you take a few moments to ask yourself
  what was the intuition that led you to that guess. You need to bring
  that intuition out of your unconscious mind and into your conscious
  mind.

  When you don't do that you are wasting your time. When you guess
  well always remember to:
  \begin{center}
    \vskip3mm{} {\sc{}\underline{Stop.}\vskip3mm{} Take a breath.  \vskip3mm{}
      \underline{Think.} \vskip3mm{}Make sure you \underline{understand.}
      \vskip3mm{} \underline{Then} move on.}  \vskip3mm{}
  \end{center}
\end{Digression}

So let's take a guess. There is no need to get really, really crazy
though. Let's make a bad guess.
How about this: $y = 1?$

Obviously that won't work because $\dfdx{y}{x}$ is equal to zero, not
one. But wait! It ``kinda'' works doesn't it? After all $y=1,$ and
$\dfdx{y}{x}=0$ are both constant functions. We just have the wrong
constant for $\dfdx{y}{x}.$ We need $\dfdx{y}{x}=1$, but we have
$\dfdx{y}{x}=0.$

Our guess has \emph{something} going for it. Let's try to tweek it a bit and
see if we can get $\dfdx{y}{x}=1.$

Actually that's easy. The derivative of $x$ is $1$ so if we we let
$y=1+x$ we get $\dfdx{y}{x} = 1.$ Again, this almost works.  We've
managed to get $\dfdx{y}{x}=1$ as desired but of course we have collateral
damage: $y$ has changed. So when we differentiate $y$ we
recover the first, constant term: $1$; but lose the second, linear
term: $x.$

Can we adjust this last guess so as to recover the second term?
Sure. Let's try $y=1+x+\frac{x^2}{2}.$ In that case
$\dfdx{y}{x} = 1+x$ and again, upon differentiating we get everything
back but the last term: $\frac{x^2}{2}.$

A few moments thought will convince you that we can pick up that term
by appending $\frac{x^3}{3\cdot2}$ to our last guess. That is, if we take
$y=1+x+\frac{x^2}{2}+\frac{x^3}{3\cdot2}$ then
$\dfdx{y}{x} = 1+x+\frac{x^2}{2}.$ But again, the last term goes
missing when we differentiate. Continuing to append
terms won't help as the following problem shows.

\begin{embeddedproblem}{}
  Compute $\dfdx{y}{x}$ and $y-\dfdx{y}{x}$ for each of the
  following\footnote{The notation $3!$ is pronounced ``three factorial''
    and means $3\cdot2\cdot1.$ Similarly
    $4!=4\cdot3\cdot2\cdot1$ and in general $n!=n(n-1)(n-2)\ldots3\cdot2\cdot1$}.
  \begin{multicols}{2}
    \begin{description}
    \item[(a)] $y= 1+x+\frac{x^2}{2!}+\frac{x^3}{3!}$
    \item[(b)] $y= 1+x+\frac{x^2}{2!}+\frac{x^3}{3!}+\frac{x^4}{4!}$
    \item[(c)]
      $y= 1+x+\frac{x^2}{2!}+\frac{x^3}{3!}+\frac{x^4}{4!}+\frac{x^5}{5!}$
    \item[(d)]
      $y= 1+x+\frac{x^2}{2!}+\frac{x^3}{3!} + \ldots +
      \frac{x^{n-1}}{(n-1)!} + \frac{x^n}{n!}$
    \end{description}
  \end{multicols}
\end{embeddedproblem}

Clearly, we've almost got something here. If we take
\begin{align*}
  y&= 1+x+\frac{x^2}{2!}+\frac{x^3}{3!} + \ldots +
     \frac{x^{999999}}{999999!} +   \frac{x^{1000000}}{1000000!} \\
  \intertext{  then}
  \dfdx{y}{x} &= 1+x+\frac{x^2}{2!}+\frac{x^3}{3!} + \ldots +
                \frac{x^{999999}}{999999!}, \\
  \intertext{and}
  y-\dfdx{y}{x} &= \frac{x^{1000000}}{1000000!}.
\end{align*}
So the  difference between $y$ and $\dfdx{y}{x}$ is that stupid
one-millionth term! Surely we can handle that somehow!

Sadly no, we can't. As we observed at the beginning of this chapter no
elementary function, or combination of functions, will solve the
IVP~\ref{eq:ExpDiffeq}. Since polynomials are elementary functions
there is no polynomial that solves the differential equation
$\dfdx{y}{x}=y.$

\begin{embeddedproblem}{}
  Show that the statement of the previous paragraph is true. That
  is, show that there is no polynomial that solves the differential
  equation: $\dfdx{y}{x}=y.$ \hint{Consider the degree of the
  polynomial and the degree of its derivative. Notice that the
  reasoning here is similar to the reasoning we used to solve
  PIC~\ref{PIC:catenary}.}
\end{embeddedproblem}

Since no polynomial will work and we already knew this, a fair
question is: Why have we taken you on this wild goose chase?

It is not completely obvious, but in fact, we \emph{almost}
have a solution here. One million is a very large number, so one
million factorial ($1000000!$) is inconceivably large. Thus for any
reasonable value of $x$ that one millionth term,
$\frac{x^{1000000}}{1000000!},$ is so incredibly close to zero that it
almost isn't really there. And if it isn't really there then
$$
y= 1+x+\frac{x^2}{2!}+\frac{x^3}{3!} + \ldots +
\frac{x^{999999}}{999999!} 
\textcolor{red}{\underbrace{+
    \frac{x^{1000000}}{1000000!}}_{\text{\tiny Not really
      here. Shhh!}}}
$$
and it's derivative
$$
\dfdx{y}{x}= 1+x+\frac{x^2}{2!}+\frac{x^3}{3!} + \ldots +
\frac{x^{999999}}{999999!} 
$$
really are the same thing. So we're .\ .\ . done. Right?

No, sadly we're not. Although that last term really is
\underline{practically} zero it is not \underline{actually} zero, no matter
how small it is. As Newton\aside{Newton wrote his scientific
  works in Latin. This is a very loose
  translation of the phrase ``\foreign{In rebus mathematicis erores quam minimi non
    sunt contemnendi.}'' from his book \pubtitle{ Quadraturam Curvarum} (On
  the Quadrature of Curves). We will have more to say about this in
  chapter~\ref{chapt:whats-wrong-with}} said, ``In mathematics the smallest of
errors must be dealt with.'' So we haven't solved our
problem in the sense of having an explicit formula for $\exp(x)$, but
we do know a great deal about it at this point. We know that $\exp(x)$
is always positive, always increasing, that it gets ever steeper, and
that near $x=0$ its graph will look very much like the graph of $
\dfdx{y}{x}= 1+x+\frac{x^2}{2!}+\frac{x^3}{3!} + \ldots +
\frac{x^{999999}}{999999!} + \frac{x^{1000000}}{1000000!}.
$
And often in real
applications an approximate solution is nearly as good as an actual
solution. This is what we have here. For any value of $x$ that is
likely to come up in any real world application
$ y= 1+x+\frac{x^2}{2!}+\frac{x^3}{3!} + \ldots +
\frac{x^{999999}}{999999!} + \frac{x^{1000000}}{1000000!}$ is an
incredibly good approximation to the solution of the \term{IVP},
$y=\dfdx{y}{x}, y(0)=1$.
So good, in fact, that we don't really have to take anywhere near
that many terms. Depending on the application involved
$y=1+x+\frac{x^2}{2!}+\frac{x^3}{3!}$ might be entirely adequate as
an approximate solution. Indeed, for some applications the
approximation $y=1+x$ is all that is really needed.

% In section~\ref{sec:eulers-method} we used Euler's Method to
% approximate the graph of the curve which solves this
% IVP. The sketch\\
% \centerline{\includegraphics*[height=1.5in,width=2in]{../Figures/exp-piecewise-linear3}}
% showed our approximation and the actual solution on the same axes.
% The sketch\\
% \centerline{\includegraphics*[height=1.5in,width=2in]{../Figures/exp-piecewise-linear4}}
% also includes the graph of $y=1+x+\frac{x^2}{2}$ in red. As you can
% plainly see  $y=1+x+\frac{x^2}{2}$ is a better approximation, at least
% on the interval $(0,1/2).$

% Still, our goal was not to approximate the solution of the
% differential equation $\dfdx{y}{x} =y,$ but to actually
% solve it. So we must keep trying.

If there is no \emph{polynomial} that solves
equation~\ref{eq:ExpDiffeq} does that mean there is no solution at
all? Certainly not. In fact, since that last term of the polynomial
was always the stumbling block it should be clear that all we need
to do is \underline{not have a last term}. This is a startling idea
but that's ok. We'll trust our intuition, but we also need to
examine it closely. What we're saying is that the solution of
equation~\ref{eq:ExpDiffeq} is:
\begin{equation}
  y= 1+x+\frac{x^2}{2!}+\frac{x^3}{3!} + \cdots\label{eq:ExpSeries1}
\end{equation}
where the dots at the end mean that the summation goes on
forever so there is no last term. 

You would expect a polynomial that doesn't end to be called and
infinite polynomial but it is not. Such an expression is called an
\term{infinite series.} Usually we will just call it a series. An
infinite series is not a polynomial by fiat. That is, we \emph{define}
polynomials to have finitely many terms. We do this preciesly so that
a series is not a polynomial.

An obvious question to ask is, "Does this infinite series even mean
anything?" Or, equivalently, "What does it mean to add up infinitely
many numbers?"  These are excellent questions which will have to be
addressed eventually. For now we won't let them trouble us. We'll just
assume that $y=1+x+\frac{x^2}{2!}+\frac{x^3}{3!}+\ldots$ makes sense,
just as we assumed that differentials make sense, and we'll come back
to this later.

\begin{embeddedproblem}{}
  Having put those deeper questions aside (for now) we can now show
  that we have found the solution of our IVP.
  \begin{enumerate}[label={\bf{}(\alph*)}]{}
\item 
  Differentiate the series
  $$
  y= 1+x+\frac{x^2}{2!}+\frac{x^3}{3!} + \ldots
  $$ 
  term-by-term to show that 
  $$
  \dfdx{y}{x}= y.
  $$ 
\item Show that $y(0)=1.$ (Yes, it really is as easy as it looks.)
\end{enumerate}
\end{embeddedproblem}

% \digress{Power Series}
% But this ``solution'' doesn't really help us understand anything at
% all. It is (for now) just empty symbol manipulation. 

% As soon as we ask ourselves what it means it is clear that we don't
% really know. For example, and most profoundly, what does it mean to
% add up infinitely many monomial terms? That is, what is an ``infinite
% polynomial.'' Does such an expression actually mean anything?

% These kinds of questions have probably never come up in your
% mathematics education before, but they will become more common from
% this point on and we will have to watch for them. Writing down an expression like 
% $$
% y= 1+x+\frac{x^2}{2!}+\frac{x^3}{3!} + \ldots
% $$ 
% is rather like asking, ``Which is bigger or a giraffe?'' It looks like
% standard English. All the parts are there: subject, verb, etc.. But
% when you put it all together it is simply meaningless. 


% If $ y= 1+x+\frac{x^2}{2!}+\frac{x^3}{3!} + \ldots $ it is to have any
% meaning we will have to decide, and agree on, what meaning we want it
% to have.
% \enddigress{}

% This approach to solving equation~\ref{eq:ExpDiffeq} will turn out to
% be very fruitful. In fact, there \underline{is} a solution and it
% \underline{can} be express as the ``infinite polynomial''
% $$
% y= 1+x+\frac{x^2}{2!}+\frac{x^3}{3!} + \cdots
% $$
% but before we can use this ``solution'' we will have to find a way to
% give meaning to the idea of ``adding up'' infinitely many terms. As
% you might imagine this will be rather tricky, so we will let it be for
% now. 

% On the other hand, the solution of equation~\ref{eq:ExpDiffeq} is of
% profound importance in all of mathematics and all of science. We will
% need to understand it as thoroughly as possible so let's try to come
% at this problem in a new way. 

% The formula 
% $$
% \dfdx{y}{x} =y
% $$ 
% says that the slope of the tangent line of our curve
% $\left(\dfdx{y}{x}\right)$ at every point is exactly equal to the $y$
% coordinate at the same point. What could such a curve look like?

% For example suppose our curve passes through the point $(0,0).$ Then
% the tangent line at $(0,0)$ will be horizontal (slope $=0.$) It is
% easiest to think of the curve as being traced out by the motion of a
% point for this example.

% We are now in a
% position to work out the essential properties of the solution of
% equation~\ref{eq:ExpDiffeq} using only the concept of the derivative
% so we we will do that now. When we are done we will have another, much
% simpler, expression for
% $ y= 1+x+\frac{x^2}{2!}+\frac{x^3}{3!} + \ldots. $


% \underline{\sc{}Fair Warning:} The development of the solution of
% equation~\ref{eq:ExpDiffeq} that you are about to read will be
% difficult to follow. You'll need to slow down and be sure you follow
% each step to fully understand what we're doing.

% Now consider the following table which gives the population of
% snaggle-toothed hornpops over a period of $10$ hours.
% $$
% \begin{array}{|ccc||ccc|}
    %     \hline
    %     \text{hours}&\text{population}&\text{rate of
                                          %                                           growth}&\text{hours}&\text{population}&\text{rate
                                                                                                                              %                                                                                                                               of
                                                                                                                              %                                                                                                                               growth}\\\hline{}
    %     0&1&&6&64&32\frac{\text{\tiny hornpops}}{\text{\tiny hour}}\\[1mm]
    %     1&2&1\frac{\text{\tiny hornpops}}{\text{\tiny hour}}&7&128&64\frac{\text{\tiny hornpops}}{\text{\tiny hour}}\\[1mm]
    %     2&4&2\frac{\text{\tiny hornpops}}{\text{\tiny hour}}&8&256&128\frac{\text{\tiny hornpops}}{\text{\tiny hour}}\\[1mm]
    %     3&8&4\frac{\text{\tiny hornpops}}{\text{\tiny hour}}&9&512&256\frac{\text{\tiny hornpops}}{\text{\tiny hour}}\\[1mm]
    %     4&16&8\frac{\text{\tiny hornpops}}{\text{\tiny hour}}&10&1024&512\frac{\text{\tiny hornpops}}{\text{\tiny hour}}\\[1mm]
    %     5&32&16\frac{\text{\tiny hornpops}}{\text{\tiny hour}}&11&2048&1024\frac{\text{\tiny hornpops}}{\text{\tiny hour}}\\[1mm]\hline
    %   \end{array}
    %     $$

    %     Clearly there is a pattern to these numbers and the pattern is easy to
    %     see. If $t$ is the number of hours elapsed then the population,
    %     $P = 2^t,$ and the rate of change of the population
    %     $\dfdx{P}{t}=2^{t-1}.$ So it appears that $\dfdx{P}{t}=2P.$ It is
    %     interesting that $P$ and $\dfdx{P}{t}$ seem to be so closely
    %     related. Let's examine that relation.

    % %     $$
    % %     \begin{array}{|ccc||ccc|}
    % %       \hline
    % %       \text{hours}&\text{population}&\text{rate of
    % %       growth}&\text{hours}&\text{population}&\text{rate
    % %       of
    % %       growth}\\\hline{}
    % %       1&2.5&1.5\frac{\text{\tiny hornpops}}{\text{\tiny hour}}&6&244.140625&97.1533224297\frac{\text{\tiny hornpops}}{\text{\tiny hour}}\\[1mm]
    % %       2&6.25&2.4871250542\frac{\text{\tiny hornpops}}{\text{\tiny hour}}&7&610.3515625&242.883306074\frac{\text{\tiny hornpops}}{\text{\tiny hour}}\\[1mm]
    % %       3&15.625&6.2178126355\frac{\text{\tiny hornpops}}{\text{\tiny hour}}&8&1525.87890625&607.208265186\frac{\text{\tiny hornpops}}{\text{\tiny hour}}\\[1mm]
    % %       4&39.0625&15.5445315888\frac{\text{\tiny hornpops}}{\text{\tiny hour}}&9&3814.697265631&1518.02066296\frac{\text{\tiny hornpops}}{\text{\tiny hour}}\\[1mm]
    % %       5&97.65625&38.8613289719\frac{\text{\tiny hornpops}}{\text{\tiny hour}}&10&9536.74316406&3795.05165741\frac{\text{\tiny hornpops}}{\text{\tiny hour}}\\[1mm]\hline
    % %     \end{array}
    % %     $$

    % %     The first thing we notice about this table is the utter futility of 
    % %     displaying so many decimals. The table is so overrun with digits
    % %     that it is nearly impossible to see any pattern to the numbers. Here
    % %     is the same chart with each number rounded to the nearest tenth.


    % %     That's better, but only slightly. 

    % %     There seems to be a relationship
    % %     between $\dfdx{P}{t}$ and $P$ similar to the relationship we noticed
    % %     in the bacteria problem earlier. However in this case it is clearly
    % %     not true that  $\dfdx{P}{t}=2P.$ 

    %     Clearly the coefficient of proportiality, $2,$ is special to this
    %     particular problem. Let's generalize things a bit by supposing instead
    %     that $\dfdx{P}{t}=kP$ for some unknown constant, $k.$

    %     Ok. Now what? It seems that we've made the problem even harder by
    %     introducing $k.$ At least we know something about the number $2,$ but
    %     $k$ is an arbitrary constant. It could be anything.

    %     Remembering that we're ultimately interested in the more general
    %     problem let's make this as simple as we can, for now, by assuming that
    %     $k=1.$ If we can solve this simpler problem maybe it will give us some
    %     insight into the general situation.


    %     \digress{}
    %     Moreover, once we have
    %     decided what is meant by an expression like 
    %     $ y=
    %     1+x+\frac{x^2}{2!}+\frac{x^3}{3!} + \ldots $ we will have 
    %     $$
    %     e^x=
    %     1+x+\frac{x^2}{2!}+\frac{x^3}{3!} + \ldots
    %     $$
    %     as well.
    %     \enddigress{}


So we have the exact solution in the form of the infinite series,
\begin{equation}
  y=1+x+\frac{x^2}{2!}+\frac{x^3}{3!}+\frac{x^4}{4!}+\ldots\label{eq:ExpSeries2}
\end{equation}
Unfortunately, we don't have the tools that allow us to work with
infinite series yet so we can only use this to approximate $\exp(x)$
by using only a finite number of terms.
    %     \begin{wrapfigure}[]{r}{3in}
    %     \captionsetup{labelformat=empty}
    %     \centerline{\includegraphics*[height=1in,width=3in]{../Figures/ExpEulerCubic}}
    %     \label{fig:ExpEulerCubic}
    %     \end{wrapfigure}
% Since we found the series solution on the curve and by
% trying to guess a polynomial solution it should be clear that 
% $$
% y=1+x+\frac{x^2}{2!}+\frac{x^3}{3!}+ \cdots + \frac{x^n}{n!}
% $$
% is also probably a good approximation. 
    %     \begin{wrapfigure}[]{l}{2in}
    %     \captionsetup{labelformat=empty}
    %     \includegraphics*[height=1in,width=2in]{../Figures/CubicExpApprox}
    %     \label{fig:CubicExpApprox}
    %     \end{wrapfigure}

% A disadvantage of Euler's Method is that a great deal of effort is
% needed to find just a few (approximate) points on our curve. But now
% that we have convincing evidence that the polynomials
% $y=1+x+\frac{x^2}{2!}+\frac{x^3}{3!}+ \cdots + \frac{x^n}{n!}$ also
% make good approximations we can use those to see what our function
% looks like farther from $(0,1).$
% Above (on the left) we've seen the
% graph of the polynomial $y=1+x+\frac{x^2}{2!}+\frac{x^3}{3!}.$ 

% Does this graph make sense to you? Stop and think about it for a
% moment.

% Our differential equation is $\dfdx{y}{x}=y.$ As we observed earlier
% this says that the slope of the curve at any point is equal to its
% $y-$coordinate. Put another way, as the graph rises it will get
% steeper.

Here is the graph of
$y=1+x+\frac{x^2}{2!}+\frac{x^3}{3!}+ \cdots + \frac{x^{50}}{50!}$
near the point $(0,1).$ You can see for yourself that this appears to
have all of the properties we've identified. The $y$ coordinate is
always positive and increasing, the graph becomes
steeper, and $y(0)=1.$\\
\centerline{\includegraphics*[height=2.5in,width=4in]{../Figures/FiftiethExp}}
But wait a minute! Have you ever seen a graph that looked like this
before?  Of course you have! When you were learning Algebra you would
have studied, for example, the exponential functions $y=2^x,$ and
$y=3^x,$ and this is exactly what they look like. The functions
$\textcolor{red}{y=1+x+\frac{x^2}{2!}+\frac{x^3}{3!}+ \cdots +
  \frac{x^{50}}{50!}}$, $\textcolor{blue}{y=2^x}$, and
$\textcolor{green}{y=3^x}$ are
all given on the same axes here:\\
\centerline{\includegraphics*[height=2.5in,width=4in]{../Figures/Exp2and3}}
As you can plainly see they are very similar.

That is, as these graphs rise, they also get steeper. So either of
$y=2^x$ or $y=3^x$ would seem to be a viable candidate for the
solution of our differential equation.  Is it possible we've had the
solution in our hands all along? Let's differentiate $y=2^x$ to see if
that's true.

  It is probably tempting to assume that we can use the Power Rule,
  giving $\dfdx{(2^x)}{x} = x2^{x-1},$ but this is \alert{wrong.}  It
  is actually pretty easy to see that this can't possibly be correct
  because when $x$ is negative then $x2^{x-1}$ is also negative. But
  the slope of $y=2^x$ is positive everywhere as you can see from its graph. In fact, none of our differentiation rules
  will give us the derivative of $y=2^x.$

So we immediately have a problem. We will have to go back to basics
and compute $\dfdx{y}{x}$ from first principles, without any of our
rules.


% \begin{wrapfigure}[]{r}{1.5in}
%   \captionsetup{labelformat=empty}
%   \centerline{\includegraphics*[height=1.5in,width=1.5in]{../Figures/DiffExp}}
%   \label{fig:DiffExp}
% \end{wrapfigure}
% Recall that we have interpreted $\d{x}$ and $\d{y}$ to be ``the
% (infinitesimal) change in $x$'' and ``the (infinitesimal) change in
% $x$,'' respectively.
Let $y(x)=2^x$ and observe that 
\begin{align*}
  y+\d{y}&=y(x+\dx{x})\\
  \intertext{Warning: $y(x+\dx{x})$ means $y$ evaluated at
  $x+\dx{x}$. This is \alert{not} multiplication. Continuing, we have}
  \d{y}&=y(x+\d{x})-y. \\
\intertext{and since $y(x+\d{x})$ is equal to $2^{x+\d{x}}$ we
  have,}
      \dx{y} &=2^{x+\d{x}}-2^x\\
       &=2^x(2^{\d{x}}-1)
\end{align*}
Thus
\begin{equation}
  \d y =2^x(2^{\d{x}}-1).  \label{eq:2^x}
\end{equation}

Similarly, if $y=3^x$ then
\begin{equation}
  \d{y} = 3^x(3^{\d{x}}-1),\label{eq:3^x}
\end{equation}

Equations~\ref{eq:2^x} and~\ref{eq:3^x} are a little scary to look
at because we're interested in computing $\dfdx{y}{x}$ but
with $\dx{x}$ in the exponent we can't do that in any obvious
fashion. We will have to be clever here.

When we originally developed the differentiation rules we thought of
$\d{x}$ as infinitely small\aside{As an infinitesimal.} so that we
could work out the procedures for finding the slope of a curve. But it
should be clear that if we take $\d{x}$ to be a {\bf{}very} small, but
\emph{finite} number, say $\d{x}\approx 0.0000001,$ we can use
equation~\ref{eq:2^x} to approximate the derivative as accurately as
we wish.

We first divide equation~\ref{eq:2^x} by $\d{x}${} to
get $ \dfdx{y}{x} =\dfdx{(2^x(2^{\d{x}}-1))}{x},$ then on the right
side of the equation set
$\d{x}\approx 0.0000001.$ This would give us
\begin{equation}
  \dfdx{y}{x}=\dfdx{(2^x)}{x}\approx 2^x\left[\frac{(2^{\d{x}}-1)}{\d
      x}\right]_{\d{x}\approx0.0000001} \approx
  2^x(0.7).\label{eq:Exp2}
\end{equation}

This isn't bad for a first try! We certainly have
$y(0)=2^0=1$, and now we see that since $y=2^x$,
$\dfdx{(2^x)}{x}\approx (0.7)2^x$. If the constant
factor was a $1$ instead of $0.7$ we'd have our formula.

Maybe $y=3^x$ will work? Actually, no. We can tell that it won't by
looking at the graphs above, but let's differentiate this
approximately just as we did with $y=2^x$.

Proceeding similarly $a=3$ we have
\begin{equation}
  \dfdx{y}{x}=\dfdx{(3^x)}{x}\approx 3^x\left[\frac{(3^{\d{x}}-1)}{\d
      x}\right]_{\d{x}\approx0.0000001} \approx
  (1.1)3^x.
  \label{eq:Exp3}
\end{equation}
That is, $\dfdx{(3^x)}{x}\approx (1.1)3^x$.
Once again the initial condition is satisfied. We have $y(0)=3^0=1$.
And once again the differential equation is {\it almost} satisified:
$\dfdx{y}{x}=3^x(1.1)$. But this time the constant, $1.1$, is too
big. This, of course, is why the green graph ($y=3^x$) is above the
red one ($y=1+x+\frac{x^2}{2!}+\frac{x^3}{3!}+ \cdots +
\frac{x^{50}}{50!}$) in our sketch above.

But there is nothing special about either ``$2$'' or ``$3$'' in the
computations above.
In general if $y=a^x$  then\aside{We restrict $a$ to positive
  values, $a>0,$ mainly because it is too hard to think about what
  this formula means if, for example, $a=-2$ and $x=\frac12.$}
\begin{equation}
  \d{y} = a^x(a^{\d{x}}-1),\label{eq:a^x}
\end{equation}

So, we have $\dfdx{(2^x)}{x}<2^x$ and $\dfdx{(3^x)}{x}>3^x$.  It
stands to reason that there must be some number between $2$ and $3$
with $\dfdx{(\text{number})^x}{x}=\text{(number)}^x$.  This number,
whatever it is, has been named $e$, mostly for historical reasons. The
function $y=e^x$ is called the \term{natural exponential}
function\aside{You can decide for yourself how ``natural'' this really
  is. Opinions vary. But whether you decide you like this name or not
  it is the name everyone uses so we're stuck with it.} and $e$ is its
base just as $2$ is the base of the exponential function\footnote{This
is \emph{an} exponential, it's just not the \emph{natural}
exponential.} $y=2^x$. Thus $y=\exp(x)=e^x$ is the
formula we've been seaching for.



We can check this by calculating
$$
\dfdx{y}{x}= e^x\left[\frac{(e^{\d{x}}-1)}{\d
    x}\right]_{\d{x}\approx0.0000001}.
$$
No doubt your calculator has an $e$ button on it. If you press that
button your calculator will respond with the numerical value of
$e$. To five decimals this is: $2.71828$. If you substitute
this calculator's value of $e$ into the formula above, you'll get
$$
\dfdx{y}{x}= e^x\left[\frac{(e^{\d{x}}-1)}{\d
    x}\right]_{\d{x}\approx0.0000001}\approx e^x(1.00000004943).
$$

We did not get $\dfdx{y}{x}$ exactly equal to $e^x$ because (1) we used
an approximation of $e$, not $e$ itself, and (2) $0.0000001$ is a very
small \alert{finite} number, not a differential.

Henceforth then, we reserve the letter $e$ to designate the base of
the natural exponential function. You might ask, ``Why not just figure
out what $e$ actually is, and use the number? Why use $e$?''  The
answer is that, much like $\pi,$  $e$ is an irrational number. This means, among
other things, that its decimal expansion never ends, so we use $e$ for
the same reason we use $\pi$.

It appears that this number $e$ between $2$ and $3$ fits the bill.
That is, $y=e^x$ satisfies the
IVP
$$
\dfdx{y}{x}=y,\ y(0)=1.
$$
\begin{mynotation}{Exponential}
  We originally named this function $\exp(x)$, and this is still a
  valid name for it, just as sqr($x$) is a valid name for the function
  $y(x)=x^2$. However, just as $y=x^2$ better represents the way we
  usually think about the squaring function (``square the input
  variable''), the notation $y=e^x$ better represents the way we
  usually think about the natural exponential function.

  For now you should think of the natural exponential as ``this funny
  number $e$, raised to the power of the input variable.''

  However, it is a curious fact that the modern definition of the
  natural exponential is \alert{not} $e^x.$ The modern definition is
  actually the infinite series we derived earlier:
  $$\exp(x)=e^x=1+x+\frac{x^2}{2!}+\frac{x^3}{3!}+ \cdots.$$

  To see why this is, ask your self what the expression $e^{1/2}$
  means.

  That's easy enough. Since $x^{1/2}=\sqrt{x}$ clearly
  $e^{1/2}=\sqrt{e}.$ It would not necessarily be easy to compute the
  square root of $e$ but we're only asking about the meaning of our
  symbols here. Since $e$ is a positive number we can take its square
  root and that is what $\sqrt{e}$ {\it{}means}, even if we can't
  compute it. In precisely the same way $e^{5/7}=\sqrt[5]{e^7}$,
  $e^{2/3}=\sqrt[3]{e^2}$ and in general if $a$ and $b$ are integers,
  $e^{a/b}=\sqrt[b]{e^a}$.

  However, suppose $q$ is an \term{irrational} number, like $\pi$.  What could the expression $e^\pi$ possibly even mean in
  that case? Since $\pi$ is not an integer it doesn't mean ``$\pi$ copies
  of $e$ multiplied together'' like say, $e^5.$ Since $q$ can't be
  represented as ratio of integers (because it is \emph{irrational}
  after all) it doesn't mean some root of $e$ raised to a power the
  way that say, $e^{a/b}$ means $\sqrt[b]{e^a}.$

  And this difficulty is not just a matter of not knowing the value of
  $e$.  It is built into the real numbers. The same issue occurs when
  we try to interpret $2^\pi$ or $3^\pi$ meaningfully because $\pi$ is
  \emph{irrational}.

  This is weird.

  Not only do we not know the value of the number $e$ we don't
   know what an expression like $e^q$ might mean\footnote{Aside: Many
    people are repelled when they encounter this sort of
    situation. Especially in mathematics. After all, mathematics is
    about having the answers to mathematical problems isn't it?

  Well, no. Actually it isn't. You're thinking of engineering.

  Mathematics is about \emph{finding} the answers to mathematical
  problems, not having them.

  As one of the greatest mathematicians
  of all time, Carl Friederick Gauss once said:\\
  \begin{center}
    \begin{minipage}[c]{0.85\linewidth}
      ``It is not the knowledge, but the act of learning, not
      possession but the act of getting there, which grants the
      greatest enjoyment. When I have clarified and exhausted a
      subject, then I turn away from it, in order to go into
      darkness again; the never-satisfied man is so strange if he
      has completed a structure, then it is not in order to dwell in
      it peacefully, but in order to begin another.''
    \end{minipage}
  \end{center}
} either, if $q$ is irratinal.

This is why the natural exponential is defined as an infinite
series. Using the series we can actually compute the numbers we need
in applications. Well, we can compute approximations to them at least.

Recall that in our first attempt to solve
equation~\ref{eq:ExpDiffeq} we concluded that the \term{infinite
  series} in equation~\ref{eq:ExpSeries2} was also a solution. Thus we
have found two solutions to the same problem: $y=e^x$ and
$y=1+x+\frac{x^2}{2!}+\frac{x^3}{3!}+\frac{x^4}{4!}+\ldots.$

If there is only one solution to the IVP (there is) then these are
both \emph{ the same solution.}  Therefore, defining $\exp(x)$ as
\begin{equation}
  \exp(x)=1+x+\frac{x^2}{2!}+\frac{x^3}{3!}+\frac{x^4}{4!}+\ldots,\label{eq:ExpSeries3}
\end{equation}
rather than $\exp(x)=e^x$  addresses all of these issues. For example
$e^\pi$ means
$$
e^\pi=
1+\pi+\frac{\pi^2}{2!}+\frac{\pi^3}{3!}+\frac{\pi^4}{4!}+\ldots.
$$
\begin{embeddedproblem}{}
  \begin{enumerate}[label={\bf{}\alph*)}]
  \item Use equation~\ref{eq:ExpSeries3} to show that $e\approx
    2.7128.$ \hint{In equation~\ref{eq:ExpSeries3} take $x=1$ and
    add up enough terms to reach $2.7128.$}
  \item Use equation~\ref{eq:ExpSeries3} to show that
    $\inverse{e}=1+(-1)+\frac{(-1)^2}{2!}+\frac{(-1)^3}{3!}+\frac{(-1)^4}{4!}+\ldots
    \approx 0.368$ \hint{In equation~\ref{eq:ExpSeries3} take
    $x=-1.$}
  \item Compute each of the following:
    \begin{enumerate}[label={\bf{}\roman*.}]
    \item $1$ (This one is easy.)
    \item $1+\frac{1}{2}$ (So is this one.)
    \item $1+(\frac12)+\frac{(\frac12)^2}{2!}$
    \item
      $1+(\frac12)+\frac{(\frac12)^2}{2!}+\frac{(\frac12)^3}{3!}$
    \item $1+(\frac12)+\frac{(\frac12)^2}{2!}+\frac{(\frac12)^3}{3!}+\frac{(\frac12)^4}{4!}$
    \end{enumerate}
    and compare them with $\sqrt{e}$ from your calculator.
  \item Find an approximation to $\sqrt[3]{e}$ that is accurate to
    $3$ decimal places.
  \end{enumerate}
\end{embeddedproblem}

\end{mynotation}

Ok, so now we have a new differentiation rule to remember:
$\dfdx{e^x}{x}=e^x$. That is, the natural exponential function
\emph{is its own derivative.}
This one is particularly easy to use but remember that just because we
have a new rule, does not mean that our old differentiation rules
are ignored.  As always all the rules work together.

\begin{myexample}
  Find the derivative of $y=e^{x}\sin(3x).$

  Using the product rule, we have
  \begin{align*}
    \dx{\left(e^{x}\sin(3x)\right)}&=e^{x}\dx{(\sin(3x))}
                                       +\sin(3x)\dx{(e^{x})}\\
                                     &=e^{x}\cos(3x)\d{(3x)} +
                                       \sin(3x)e^{x}\dx{x}\\
                                     &=e^{x}\cos(3x)3\dx{x} +
                                       \sin(3x)e^{x}\dx{x}\\
    \intertext{so that}
    \dfdx{\left(e^{x}\sin(3x)\right)}{x}&=e^{x}\left[3\cos(3x) +
                                            \sin(3x)\right]\\ 
  \end{align*}

\end{myexample}

\begin{ProficiencyDrill}
  Compute $\dfdx{y}{t}$ for each of the following. \hint{Make them
    easier on your eyes.}
  \begin{enumerate}[label={\bf{}(\alph*)}]
      \begin{multicols}{2}
      \item $y=e^{2t}$
      \item $y=e^{3t}$
      \item $y=e^{\pi t}$
      \item $y=e^{et}$
      \end{multicols}
    \item Find the pattern in parts (a) through (d) and use it to
      compute $\dfdx{y}{t}$ if $y=e^{\alpha{}t}.$ Assume $\alpha{}$ is
      a constant.
        \begin{multicols}{2}
        \item $y=te^t$
        \item $y=(t+1)e^t$
        \item $y=(t^2+t+1)e^t$
        \item $y=(t^3+t^2+t+1)e^t$
        \end{multicols}
      \item Find the pattern in parts (f) through (i) and use it to
        compute $\dfdx{y}{t}$ if $$y=(t^4+t^3+t^2+t+1)e^t.$$
      \item Can you extend this pattern to find $\dfdx{y}{t}$ if
        $y=P(t)e^t$ where $P(t)$ is \underline{any} polynomial? 
      \item Can you extend this pattern to find $\dfdx{y}{t}$ if
        $y=f(t)e^t$ where $f(t)$ is {any} function? 
    \end{enumerate}
  \end{ProficiencyDrill}

  
% This seems to suggest that there is some number between $2$ and $3$  which provides a
% ``natural'' exponential function which satisfies the differential
% equation $\dfdx{y}{x}=y.$  How could we figure out what that number
% is?

% First let's give our hypothetical number a name. The traditional
% name is $e.$ If everything we've done so far is valid then
% the solution of $\dfdx{y}{x}=y$ is 
% $$
% y=e^x
% $$
% where $e$ is some, currently unknown, number.

    %     Once again, we've \emph{almost} found the solution to our differential
    %     equation: $\dfdx{y}{x}=y.$

    %     This is clear because when  $y=2^x,$ then its derivative $\dfdx{y}{x}$ is a little \emph{less} than $2^x$
    %     (since $0.7\ldots$ is a little less than $1$). That is, when $a=2,$
    %     $\dfdx{y}{x}<y.$

    %     Similarly when $y=3^x,$ $\dfdx{y}{x}$ is a little \emph{more} than $3^x$
    %     (since $1.1\ldots$ is a little more than $1$). That is, when $a=3,$
    %     $\dfdx{y}{x}>y.$

    %     This seems to suggest that there is some number between $2$ and $3,$
    %     -- let's call\footnote{Cuz that is what everyone else in the world
    %     calls it.} it $e$ -- where $\dfdx{y}{x}=y.$ That is $y=e^x$ is
    %     actually a \emph{solution} of our differential
    %     equation~\ref{eq::ExpDiffeq2}.

% The exponential function $y=e^x$ is Icalled the ``natural'' exponential
% function. It is natural in the sense that it is its own
% derivative. That
% is, it satisfies the differential equation $\dfdx{y}{x}=y.$ Notice that we do not know the value of $e,$ only that it is between $2$ and $3.$ How could
% we figure out what the value of the number $e$ is?


% Did you notice that we've slipped something past you here?

  % Recall that we saw in equation~\ref{eq:Exp2} that the derivative of
  % the  function
  % $y=2^x$ is (approximately)
  % $$
  % \dfdx{y}{x}=  2^x(0.7).
  % $$
  % It should be clear that it also (approximately) solves the IVP:
  % $$
  % \dfdx{y}{x}=a2^x,\text{  }y(0)=0.7
  % $$
  % IVP~(\ref{eq:ExpDiffeq2} consist of two parts, a
  % differential equation and an initial condition. What happens if the
  % initial condition is different?
  
  % Can we solve the IVP 


% No, we don't. In fact, it is clear that $y(0)=e^0=1,$ since \emph{any}
% number raised to the zeroth power is equal to\aside{Is that really true? What about zero to the zeroth power. Hmm\ .\ .\ .} $1.$

% But don't miss this important point: We were just lucky that this worked
% out. Had our initial value been anything else then $y=e^x$ would have
% solved the differential equation, but not the initial condition, and
% hence it would not have solved the IVP.

\begin{embeddedproblem}{}
  Use\notetoself{This doesn't belong here -- ecb -- May 29, 2020} equations~\ref{eq:Exp2} and~\ref{eq:Exp3} to find an IVP for
  which each of the following is the solution:
  \begin{enumerate}[label={\bf{}(\alph*)}]
    \begin{multicols}{3}
  \item $y=2^x$
  \item $y=3^x$
  \item $y=a^x$
    (We need to  assume that $a>0.$ Why?)
    \end{multicols}
  \end{enumerate}
\end{embeddedproblem}

% Because  $e$ is a number, just like $2$ or $3,$ we
% can think of the expression $e^x$ in exactly the same way we
% think about the expression $2^x.$ That is, if
% $2^3=2\cdot2\cdot2$ then clearly $e^3=e\cdot e\cdot e$
% Similarly, just as $3^{1/2}=\sqrt{3}$ it is also true that
% $e^{1/2}=\sqrt{e},$ and in general if $m$ is any
% \term{rational} number (by which we mean that $m=a/b$ where
% $a$ and $b$ are both integers.) then
% $e^m=e^{a/b} = \sqrt[b]{e^a},$ just as
% $2^{a/b} = \sqrt[b]{2^a}.$




%     %     but which number?
%     %     What is its value? 

%     %     First let's give our hypothetical number a name. The traditional
%     %     name is $e.$ If everything we've done so far is valid the
%     %     the solution of $\dfdx{y}{x}=y$ is 
%     %     $$
%     %     y=e^x
%     %     $$
%     %     where $e$ is some, currently unknown, number.

%     %     But before we dismiss these as trivial let's look at the
%     %     \emph{meaning} of the symbols, not just their relationships.

%     %     If $x$ is moving and has some value, the $x+\d x$ is its value ``just
%     %     a little farther along.'' Thus equation~\ref{infx} says that the
%     %     differential of $x,$  $\d x,$ is the difference between the value of
%     %     $x$ a little farther along and its value now. Again, this is pretty
%     %     obvious.

%     %     Equation~\ref{infy} says the same thing about $y,$ but this time it is
%     %     not so easy to find the value of $y,$ $y+\d y,$ ``a little farther
%     %     along.'' This is because $y$ \alert{depends} on $x$ in a
%     %     non-elementary way.


%     %     what it 
%     %     if $y_2$ has some value then $x+\d{x}$ is it's
%     %     value \emph{after} it has been incremented by $\d{x}.$ And so,
%     %     $$
%     %     \d{x} = (x+\d{x}) - x.
%     %     $$



%     %     If we increment $x$  by the infinitesimal amount $\d{x},$ we get
%     %     \begin{align*}
%     %     \d{2^x} &= 2^{x+\d{x}}-2^x \\
%     %                       &= 2^x\left(2^{\d{x}}-1\right)\\
%     %     \intertext{and dividing both sides by $\d{x}$ gives the derivative}
%     %     \dfdx{(2^x)}{x} &= 2^x\left(\frac{2^{\d{x}}-1}{\d{x}}\right).
%                             %                             \intertext{In an exactly analogous fashion we get}
%                             %                             \dfdx{(3^x)}{x} &= 3^x\left(\frac{3^{\d{x}}-1}{\d{x}}\right).
%                                                                             %   \end{align*}

%                                                                             %                                                                             Of course we ordinarily think of $\d{x}$ as infinitely small, but we
%                                                                             %                                                                             can approximate the derivative by taking it to be a very small
%                                                                             %                                                                             \underline{finite} number, say $\d{x}\approx 0.0000001.$ This would
%                                                                             %                                                                             give us
%                                                                             %                                                                             $$
%                                                                             %                                                                             \dfdx{(2^x)}{x} \approx
%                                                                             %                                                                             2^x\left(\frac{2^{0.0000001}-1}{0.0000001}\right) = 2^x(0.693\ldots)
%                                                                             %                                                                             $$
%                                                                             %                                                                             and
%                                                                             %                                                                             $$
%                                                                             %                                                                             \dfdx{(3^x)}{x} \approx
%                                                                             %                                                                             3^x\left(\frac{3^{0.0000001}-1}{0.0000001}\right) = 3^x(1.098\ldots).
%                                                                             %                                                                             $$

%                                                                             %                                                                             Observe that when $y=2^x,$ $\dfdx{y}{x}$ is a little \emph{less} than $2^x$
%                                                                             %                                                                             (since $0.693\ldots$ is a little less than $1$).
%                                                                             %                                                                             and when $y=3^x,$ $\dfdx{y}{x}$ is a little \emph{more} than $3^x$
%                                                                             %                                                                             (since $1.098\ldots$ is a little more than $1$).
%                                                                             %                                                                             This seems to suggest that there is some number between $2$ and $3$  which provides a
%                                                                             %                                                                             ``natural'' exponential function which satisfies the differential
%                                                                             %                                                                             equation $\dfdx{y}{x}=y.$  How could we figure out what that number
%                                                                             %                                                                             is?

%                                                                             %                                                                             First let's give our hypothetical number a name. The traditional
%                                                                             %                                                                             name is $e.$ If everything we've done so far is valid the
%                                                                             %                                                                             the solution of $\dfdx{y}{x}=y$ is 
%                                                                             %                                                                             $$
%                                                                             %                                                                             y=e^x
%                                                                             %                                                                             $$
%                                                                             %                                                                             where $e$ is some, currently unknown, number.

% \begin{myexample}
%   Of course, as usual our reasoning here is a bit flaky. We argued,
%   just from the shape of the graphs, that the curves $y=2^x$ and
%   $y=3^x$ must be related to the solution of the differential
%   equation~\ref{eq:ExpDiffeq}. At best this is very vague
%   reasoning. The graph of the curve $y=e^{\frac{x^2}{2}}$ also looks
%   superficially like the graphs of $y=2^x$ and
%   $y=3^x$ so we could just as easily (and just as non-rigorously) have
%   argued  that this must be the solution of our differential
%   equation.

%   Except that it isn't.

%   It is instead the solution of the differential equation
%   $$
%   \dfdx{y}{x} = xy.
%   $$
%   \begin{embeddedproblem}{}
%     Verify that statement. That is, multiply $x$ times $y,$ and take
%     the derivative of $y=e^\frac{x^2}{2}$ and verify that they are
%     equal.
%   \end{embeddedproblem}
% \end{myexample}

%                                                                             %                                                                             We've also seen that $e^x=1+x+\frac{x^2}{2!}+\frac{x^3}{3!} + \ldots.$

%                                                                             %                                                                             Notice that we can obtain the base of $2^x$ or $3^x$ (or \emph{any}
%                                                                             %                                                                             exponential function, $a^x$) simply by taking
%                                                                             %                                                                             $x=1.$ That is, $2^1=2,$ $3^1=3,$ and $a^1=a.$ Of course this is not
%                                                                             %                                                                             very interesting for those functions but in the case of $e^x$ we can
%                                                                             %                                                                             find the numerical value of $e$ in the same way, by setting $x$ equal
%                                                                             %                                                                             to one: $e^1 = 1+1+\frac{1^2}{2!}+\frac{1^3}{3!} + \ldots.$ Of course
%                                                                             %                                                                             since there are infinitely many terms we can't add them \emph{all} up,
%                                                                             %                                                                             but if we take enough of them we can approximate $e$ to any degree of
%                                                                             %                                                                             accuracy. It turns out that\footnote{We're really just
%                                                                             %                                                                             showing off by giving this many decimals. For any practical application $e\approx 2.718$ is more
%                                                                             %                                                                             than enough.} 
%                                                                             %                                                                             $$
%                                                                             %                                                                             e=2.718281828459045\ldots.
%                                                                             %                                                                             $$

%                                                                             %                                                                             Notice that $2<e<3$ as we expected.

%                                                                             %                                                                             Let's proceed on that assumption\footnote{Obviously this is actually
%                                                                             %                                                                             true or we wouldn't have spent this much time on it.} for now,
%                                                                             %                                                                             recognizing that this is only a conjecture for now. That is, we make
%                                                                             %                                                                             the following conjecture.

%                                                                             %                                                                             \begin{myconjecture}
%                                                                             %                                                                             There is a real number we'll denote by $e,$ with the property that
%                                                                             %                                                                             $e^x$ satisfies equation~\ref{eq:ExpDiffeq}. That is
%                                                                             %                                                                             $$
%                                                                             %                                                                             \dfdx{(e^x)}{x} = e^x.
%                                                                             %                                                                             $$
%                                                                             %                                                                             \end{myconjecture}

%                                                                             %                                                                             Depending on your mathematical background you may have seen this
%                                                                             %                                                                             number, $e,$ before and you may already know that it is approximately
%                                                                             %                                                                             $2.7128.$ This is true but it isn't especially useful information for
%                                                                             %                                                                             our purposes right now.

%                                                                             %                                                                             What is useful is that we now have strong evidence that the solution
%                                                                             %                                                                             of equation~\ref{eq:ExpDiffeq} is $y=e^x,$  and that this can be
%                                                                             %                                                                             interpreted as an exponential function just like $10^x$ or $2^x.$
%                                                                             %                                                                             Because of this the expression $e^x$ has been named the ``natural
%                                                                             %                                                                             exponential.''

% As we observed earlier the function $e^x$ is called the
% \underline{natural exponential function} even though at first glance
% hardly anything about it seems natural.  One could say that satisfying
% $dy/dx=y$ makes it natural, but there is more to it than that, as we
% will see later\footnote{\textcolor{red}{We need to come back to this
%     and show that all exponential can be expressed in terms of the
%     natural. Also when the logarithm is seen as the antiderivative of
%     1/t.}}.

% First, consider that by the Chain Rule, if $x=kt$ for some constant
% $k,$ then
% $\dfdx{y}{t}=\dfdx{e^{kt}}{t}=\dfdx{(e^x)}{x}\dfdx{x}{t}=e^x\cdot{}k=ky.$


\begin{embeddedproblem}{}
  In this problem we will (finally!) build the solution of the Hanging
  Chain problem. 
  \begin{enumerate}[label={\bf{}(\alph*)}]

  \item Show that
    $1+\left(\frac{e^x+e^{-x}}{2}\right)^2=\left(\frac{e^x-e^{-x}}{2}\right)^2.$\\
    \hint{It might help to make the substitution $e^x=a$ just to
      make things ``easier on the eyes.'' That makes
      $e^{-x}=\frac1a.$}
  \item Show that $
    \dfdx{
      \left(\frac{1}{2}(e^x+ e^{-x})\right)}{x}=
    \frac{1}{2}(e^x- e^{-x}) $ and that  $\dfdx{
      \left(\frac{1}{2}(e^x- e^{-x})\right)}{x}=
    \frac{1}{2}(e^x+ e^{-x})$
  \item Show that the curve $y=\frac12\left(e^x+e^{-x}\right)$ satisfies the
    differential equation:
    $$
    \dfdxn{y}{x}{2}=\sqrt{1+\left(\dfdx{y}{x}\right)^2}
    $$ 
  \item Now show (at last!) that the curve
    \begin{equation}
      y=\frac{H}{w}
      \left(\frac{e^{\frac{wx}{H}}+e^{-\frac{wx}{H}}}{2}\right)\label{eq:CatSol}
    \end{equation}

    really does satisfies the equation of the catenary
    $$
    \dfdxn{y}{x}{2}=\frac{w}{H} \sqrt{1+\left(\dfdx{y}{x}\right)^2}
    $$
    where $w$ is the weight density of the chain, and $H$ is the
    constant (magnitude of the) horizontal tension.
  \item   Assume $w=1,$ and  $H=1$ and graph the curve given by
    equation~\ref{eq:CatSol}.  Does it look like a hanging chain?
    What happens to the graph if we use $w=1, H=2$ or $w=2, H=1.$  Does this
    make sense physically?  Why or why not?
  \end{enumerate}

\end{embeddedproblem}

  \begin{embeddedproblem}{}
    Compute each of the following:
    \begin{enumerate}[label={\bf{}(\alph*)}]
      \begin{multicols}{4}
      \item $\dx(e^{x^2+\sin(x)})$
      \item $\dx(e^{x}\cos(3x)$
      \item $\dx
        \left(
          \frac{3\sqrt{x}+1}{e^x}
        \right)$
      \item $\dx(\tan(e^x))$

      \end{multicols}
    \end{enumerate}{}
  \end{embeddedproblem}

  \begin{embeddedproblem}{}
    For each of the    following, find an equation relating $\dx{x},$
    $\dx{y},$  $\dx{z}.$ Use this to compute $\dfdx{y}{x}$  and
    $\dfdx{y}{t}.$
    \begin{enumerate}[label={\bf{}(\alph*)}]{}
      \begin{multicols}{2}
      \item $e^y=x^2+3x-2$
      \item $\sin(xy)=(e^x)^y$
      \item $\tan(x+z)=e^ye^x$
      \item $e^{x^2+y^2+z^2}=2$
      \end{multicols}
    \end{enumerate}
  \end{embeddedproblem}

  \begin{embeddedproblem}{}
    Use Newton's Method to approximate the coordinates of the
    intersection point of the curves $y = -x$ and $y = e^x$.\\
    \centerline{\includegraphics*[height=1.5in,width=2in]{../Figures/xex-Intersect}}
  \end{embeddedproblem}

  \begin{embeddedproblem}{}
    Show that $y(t)=y_0e^{.05t}$ grows continually at a rate of
    $5\%$. What can you say about $y_0e^{rt}$ in general? We will
    explore this in more depth later.) 
  \end{embeddedproblem}

  \begin{embeddedproblem}{}
    Notice that $y=e^{t+a}$ satisfies $\dfdx{y}{t}=y$ for any constant
    $a$ as well. How do you make peace with this?
  \end{embeddedproblem}

  \begin{embeddedproblem}{}
    Show that the graph of $y=e^{ax}$ is increasing when $a>0$ and
    decreasing when $a<0$. Does this make sense to you?
  \end{embeddedproblem}


Recall in Problem~\ref{prob:HarmonicOscillator} we saw that the
differential equation that models the motion of a mass $m$   on a
spring oscillating without any resistance is:
$$
m\dfdxn{y}{t}{2}= -ky,
$$
where the constant $k$ is the spring constant.  This was an example of
a simple harmonic oscillator. It is ``simple'' because without
resistance, the mass would oscillate forever which we know won't
happen. 
This is illustrated by the fact that its solution is given by
$$
y=A\sin\sqrt{\frac{k}{m}}t+B\cos\sqrt{\frac{k}{m}}t
$$
for any constants $A$ and $B.$

But in the real world the spring's oscillations will get shorter and
shorter, eventually stopping altogether. How could we modify this
model in order to capture the effects of some sort of resistance, say
air resistance or resistance from some fluid as in a car's shock
absorber?

A relatively simple way to model this resistance is to say that the
resistance force is proportional to the speed of the object.  This
sounds imposing when we say it like that, but you have probably
experienced this yourself in a moving vehicle.  If you stick your hand
out the window, you notice that as the car moves faster, there is more
wind (force) pushing on your hand.  This simple model breaks down at
high speeds where factors such as turbulence affect the resistance,
but experimentally, this assumption works well for our spring model.
Recall that the term $-ky$ was the force due to the spring (Hooke's
Law) and has the opposite sign of $y.$ The explanation is that if
$y>0$ then the mass is above the equilibrium point and the spring
pushes down and if $y<0$ then the mass is below the equilibrium point
and the spring pulls upward.

By Newton's Second Law, Force equals Mass ($m$) times Acceleration
($\dfdxn{y}{t}{2}$), so the force acting on the spring is
$m\dfdxn{y}{t}{2}$. Adding the resistance simply means that we need to
add in a (negative) force which is proportional to (a constant
multiple of) the velocity, $\dfdx{y}{t}$. So we add the term
$-b\dfdx{y}{t}$ to the right hand side of the equation, where $b>0$ is
a constant.  This represents the resistance force which is
proportional to velocity of the object.

Thus our model for an oscillation with resistance becomes
$$
m\dfdxn{y}{t}{2}=-b\dfdx{y}{t}-ky
$$
or
\begin{equation}
\dfdxn{y}{t}{2}+\frac{b}{m}\dfdx{y}{t}+\frac{k}{m}y = 0.\label{eq:DampedOscillator}
\end{equation}

% \begin{embeddedproblem}{}
%   Explain why the constant of proportionality is negative in the
%   resistance.
% \end{embeddedproblem}


\begin{embeddedproblem}{}
  To simplfy the computations a bit, for this problem we'l take
  $b=k=m$ so that the  equation~\ref{eq:DampedOscillator} becomes:
  \begin{equation}
    \dfdxn{y}{t}{2}+\dfdx{y}{t}+y=0.\label{eq:DampedProblem1}
  \end{equation}
  \begin{enumerate}[label={\bf{}(\alph*)}]{}
  \item Show that: $\displaystyle y=e^{\frac{-t}{2}}
    \left[
      A\sin\left(\frac{\sqrt{3}}{2}t\right)+B\cos\left(\frac{\sqrt{3}}{2}t\right)
    \right]$ satisfies the differential equation
    $$
    \dfdxn{y}{t}{2}+\dfdx{y}{t}+y=0
    $$
    for any constants $A$ and $B$.
  \item Let $A=1, B=0$ and plot this solution for $0\le t\le 15$
    and $-0.5\le y \le 0.5.$ Does this seem to model a damped
    oscillation?
  \end{enumerate}
\end{embeddedproblem}

While it is an easy matter to verify  a solution as in the previous
problem,  it is quite another to find such a solution in the first place.
That is much harder to do.

Leonhard Euler $(1707-1783)$ solved the differential
equation\ref{eq:DampedProblem1} by guessing. Here's how he did
it. Undoubtedly Euler made several guesses that didn't help, but
eventually he hit on the guess:
$$
y=e^{ct}.
$$
Now, we already know that
$y=e^{\frac{-t}{2}} \left[
  A\sin\left(\frac{\sqrt{3}}{2}t\right)+B\cos\left(\frac{\sqrt{3}}{2}t\right)
\right]$ is a solution of our IVP so this looks like a crazy guess
because there doesn't seem to be any way to get from an exponential
function to a trigonometric function. But Euler was a genius so he
didn't let that bother him. He gave himself a little ``wiggle room''
by leaving the constant $c$ unspecified.

Sure enough, when he put his guess into
equation~\ref{eq:DampedOscillator} he found that to compute $c$ he had
to take the square root of a negative number, which doesn't seem to
make any sense. At this point most people would
assume that they'd made some mistake and would give up, or start
again. But Euler was one of the all-time great mathematicians.  In any
list of stellar mathematicians, he would stand out among the best.
Indeed, Pierre-Simon Laplace $(1749-1827)$, an outstanding
mathematician in his own right, was once asked by a novice how best to
learn Calculus. Laplace told him ``Read Euler, read Euler, he is the
master of us all.''

When Euler encountered the expression $\sqrt{-1}$ he just wrote it
down, assumed that it represented some constant, just like $1,$ or $\pi$,
but otherwise he didn't concern himself about it. He just proceeded on
with his calculation.

\begin{embeddedproblem}{}
  \begin{enumerate}[label={\bf{}(\alph*)}]
      \item   Show that
        \begin{align}
\label{eq:DampedProblem2}          y_1&=e^{(\frac{-1}{2}+\frac{\sqrt{-1}\sqrt{3}}{2})t}\text{ and}\\
\nonumber          y_2&=e^{(\frac{-1}{2}-\frac{\sqrt{-1}\sqrt{3}}{2})t}
        \end{align}
        both satisfy  the differential
        equation~\ref{eq:DampedProblem1}. (Note the presence of and
        location of $\sqrt{-1}.$
        \item Now show every combination of $y_1$ and $y_2$ of the
          form, $y=Ay_1+By_2$ where  $A$ and $B$  constants also
          satisfies equation~\ref{eq:DampedProblem1}.
  \end{enumerate}
\end{embeddedproblem} 

Since
$$
y=Ae^{\left(\frac{-1}{2}+\frac{\sqrt{-1}\sqrt{3}}{2}\right)t}+Be^{\left(\frac{-1}{2}-\frac{\sqrt{-1}\sqrt{3}}{2}\right)t}
= e^{\left(\frac{-1}{2}\right)t}\left[Ae^{\left(\frac{\sqrt{-1}\sqrt{3}}{2}\right)t}+Be^{\left(\frac{-\sqrt{-1}\sqrt{3}}{2}\right)t}\right]
$$
we  can see where the exponential part of  $e^{\frac{-1}{2}t}$ part
equation~\ref{eq:DampedProblem2} comes from, but what about the sine
and cosine? And those imaginary numbers? How do they fit in? This is
very strange.

Suppose we go back to the simplest such differential equation
$$ \dfdxn{y}{t}{2}=-y.
$$
Euler knew, as we know, that the solution of this looks like
$$
y=A\cos(t)+B\sin(t)
$$
for any choice of constants $A$ $B.$ But instead of going with that
known solution, Euler tried the same kind of guess as before.  Suppose
the solution looked like $y=e^{ct}$ for some constant $c.$ What must
that constant be?

\begin{embeddedproblem}{}
  Substitute $y=e^{ct}$ into $\dfdxn{y}{t}{2}=-y$ show that this
  function satisfies the differential equation when $c=\sqrt{-1}$
  or when $c=-\sqrt{-1}.$
\end{embeddedproblem}

You should have obtained that either one of $e^{\sqrt{-1}t}$ or $e^{-\sqrt{-1}t}$
should satisfy the above differential equation.  Let's focus on the
first.  According to what we know
$$
e^{\sqrt{-1}t}=A\cos(t)+B\sin(t)
$$
for some judicious choice of $A$ and $B.$

\begin{embeddedproblem}{}
  \begin{enumerate}[label={\bf{}(\alph*)}]{}
  \item Show that if $e^{\sqrt{-1}t}=A\cos(t)+B\sin(t),$ then $A=1$ and $B=\sqrt{-1}.$
    \hint{To get $A,$ substitute $t=0.$ To get $B,$ differentiate
      first.}
  \item Do the same to obtain $e^{-\sqrt{-1}t}=\cos(t)-\sqrt{-1}\sin(t).$
  \item Show that
    \begin{align*}
      \frac{e^{\sqrt{-1}t}+e^{-\sqrt{-1}t}}{2} &= \cos(t) \text{ and} \\
      \frac{e^{\sqrt{-1}t}-e^{-\sqrt{-1}t}}{2\sqrt{-1}}&= \sin(t).
    \end{align*}
  \end{enumerate}
\end{embeddedproblem}

This formula, $e^{\sqrt{-1}t}=\cos(t)+ \sqrt{-1}\sin(t),$ is known as Euler's formula
(for complex numbers) and is central to understanding the geometry of
complex numbers.  It is also of great importance in applied
mathematics and physics.  For Euler, it was a way to translate
trigonometric functions to exponential functions and back.  For us, it
is the key to our dampened oscillation problem.

\begin{embeddedproblem}{}
  Let's go back to our solution
  $y=
  e^{\left(\frac{-1}{2}\right)t}\left[Ae^{\left(\frac{\sqrt{-1}\sqrt{3}}{2}\right)t}+Be^{\left(\frac{-\sqrt{-1}\sqrt{3}}{2}\right)t}\right]$
  of the equation $\dfdxn{y}{t}{2}+\dfdx{y}{t}+y=0.$ Use Euler's
  formula to transform this solution into one of the form
  $$y=e^{\frac{-t}{2}} \left[
    A\sin\left(\frac{\sqrt{3}}{2}t\right)+B\cos\left(\frac{\sqrt{3}}{2}t\right)
  \right].$$ Your constants $A$ and $B$ will involve $\sqrt{-1}$, but don't let
  that bother you.  If you are dubious on this, then consider the
  initial conditions $y(0)=0$ and
  $\left.\dfdx{y}{t}\right|_{t=0}=\frac{\sqrt{3}}{2}$ and use these to
  determine $A$ and $B$.
  $$
  y=ae^{
    \left(
      \frac{-1}{2}+\frac{\sqrt{-1}\sqrt{3}}{2}
    \right)t
  }+be^{
    \left(
      \frac{-1}{2}-\frac{\sqrt{-1}\sqrt{3}}{2}
    \right)t
  }.
  $$
  Use these values of $a$ and $b$ to determine $A$ and $B.$ Does this
  answer satisfy the initial conditions?
\end{embeddedproblem}

\begin{myexample}[Compound Interest and Exponential Functions]
\label{example:CompoundInterest}
  The natural exponential function also arises naturally in financial
  mathematics. 
  Suppose we invest money in a bond which nominally pays 5\% annually,
  compounded quarterly.  The effective yield would be 5.09\%.  What
  does any of this mean?
  
\begin{wraptable}[]{r}{4.5in}
  \begin{tabular}[h]{|c|c|}
    \multicolumn{2}{c}{\textcolor{blue}{compounding $\$1$ semiannually}}\\\hline{}
        Time $t$ in years &Amount investment is worth in
                            dollars\bigstrut{}\\\hline
        $0$&\$1\bigstrut{}\\
        $1/2$&$ \$1+\$1\left(\frac{0.05}{2}\right)=\$1\left(1+\frac{0.05}{2}\right)$\bigstrut{}\\
        $1$&$  \$1(1+0.05/2)+\$1(1+0.05/2)(0.05/2)=\$1\left(1+\frac{0.05}{2}\right)^2$\bigstrut{}\\\hline
      \end{tabular}
      \caption{}
\label{tab:CompInt6mo}
\end{wraptable}
To start, the effective yield represents the actual amount of money
  earned at the end of one year.  For example, if the investment was
  not compounded at all, then at the end on one year, every dollar
  invested would yield $\$1.05$ in return, so the yield really is 5\%
  annually.  If the investment is compounded semiannually, then this
  means that half of interested earned in paid out halfway through the
  year.  This money is then re-invested for the next half year.  This
  is summarized  table~\ref{tab:CompInt6mo}.
%   \begin{table}[h]
%           \label{tab:CompInt6mo}
% \centering
%       \begin{tabular}[h]{|c|c|}\hline
%         Time $t$ in years &Amount investment is worth in
%                             dollars\bigstrut{}\\\hline
%         $0$&\$1\bigstrut{}\\\hline
%         $1/2$&$ \$1+\$1\left(\frac{0.05}{2}\right)=\$1\left(1+\frac{0.05}{2}\right)$\bigstrut{}\\\hline
%         $1$&$  \$1(1+0.05/2)+\$1(1+0.05/2)(0.05/2)=\$1\left(1+\frac{0.05}{2}\right)^2$\bigstrut{}\\\hline
%       \end{tabular}
%       \caption{\textcolor{blue}{compounding $\$1$ semiannually}}
%     \end{table}

          
%    $$
    % \$1(1+0.05/2)+\$1(1+0.05/2)(0.05/2)
    % $$
\begin{wraptable}[]{l}{3.5in}
  \begin{tabular}[h]{|c|c|}
    \multicolumn{2}{c}{\textcolor{blue}{Compounding $\$1$ three times annually}}\\\hline
      Time in years & Amount investment is worth in dollars\bigstrut{}\\\hline{}
      $0$ & $\$1$\bigstrut{}\\
      $1/3$ & $\$1\left(1+\frac{0.05}{3}\right)^1\approx1.0167$\bigstrut\\
      $2/3$ & $\$1\left(1+\frac{0.05}{3}\right)^2\approx1.0336$\bigstrut\\
      $1$   & $\$1\left(1+\frac{0.05}{3}\right)^3\approx1.0508$\bigstrut\\\hline
    \end{tabular}
  \caption{}
\label{tab:CompInt4mo}
\end{wraptable}

Tables~\ref{tab:CompInt4mo} and~\ref{tab:CompInt3mo} show, in abbreviated form, the
computations needed if the investment is compounded three times in a
year. Notice how these three tables differ from and are  similar to each
other.
% \begin{table}[h]\centering
%     \begin{tabular}[h!]{|c|c|}\hline
%       Time in years & Amount investment is worth in dollars\bigstrut{}\\\hline{}
%       $0$ & $\$1$\bigstrut{}\\\hline
%       $1/3$ & $\$1\left(1+\frac{0.05}{3}\right)^1\approx1.0167$\bigstrut\\\hline
%       $2/3$ & $\$1\left(1+\frac{0.05}{3}\right)^2\approx1.0336$\bigstrut\\\hline
%       $1$   & $\$1\left(1+\frac{0.05}{3}\right)^3\approx1.0508$\bigstrut\\\hline
%     \end{tabular}
%   \caption{\textcolor{blue}{Compounding $\$1$ three times annually}}
% \end{table}


Compounded quarterly means that the interest is re-invested four
times.  Our table looks like this.

% \begin{table}[h]\centering
%     \begin{tabular}[h]{|c|c|}\hline
%       Time in years & Amount investment is worth in dollars\bigstrut{}\\\hline{}
%       $0$   & $\$1$\bigstrut{}\\\hline
%       $1/4$ & $\$1\left(1+\frac{0.05}{4}\right)^1\approx1.0125$\bigstrut\\\hline
%       $1/2$ & $\$1\left(1+\frac{0.05}{4}\right)^2\approx1.02552$\bigstrut\\\hline
%       $3/4$ & $\$1\left(1+\frac{0.05}{4}\right)^3\approx1.038$\bigstrut\\\hline
%       $1$   & $\$1\left(1+\frac{0.05}{4}\right)^4\approx1.0509$\bigstrut\\\hline
%     \end{tabular}
%   \caption{\textcolor{blue}{Compounding $\$1$ quarterly}}
% \end{table}

The effective yield comes from the last entry in each table.  So if we
compound quarterly it is the difference between the amount the
investment is worth at the end of the year and what it is worth at the
beginning of the year: $1.0509-1=0.0509=5.09\%. $

\begin{wraptable}[]{r}{3.5in}
  \begin{tabular}[h]{|c|c|}
    \multicolumn{2}{c}{\textcolor{blue}{Compounding $\$1$ quarterly}}\\\hline
      Time in years & Amount investment is worth in dollars\bigstrut{}\\\hline{}
      $0$   & $\$1$\bigstrut{}\\
      $1/4$ & $\$1\left(1+\frac{0.05}{4}\right)^1\approx1.0125$\bigstrut\\
      $1/2$ & $\$1\left(1+\frac{0.05}{4}\right)^2\approx1.02552$\bigstrut\\
      $3/4$ & $\$1\left(1+\frac{0.05}{4}\right)^3\approx1.038$\bigstrut\\
      $1$   & $\$1\left(1+\frac{0.05}{4}\right)^4\approx1.0509$\bigstrut\\\hline
    \end{tabular}
  \caption{}
\label{tab:CompInt3mo}
\end{wraptable}
Following the same reasoning, if the investment was compounded daily,
then the effective yield would be
$\left(1+\frac{0.05}{365}\right)^{365}-1\approx0.0513 = 5.13\%$.  In
general, if we compound $n$ times in a year, then the return on an
investment of one dollar at the end of one year would be
$\left(1+\frac{0.05}{n}\right)^{n}$ dollars, after $2$ years
$\left(1+\frac{0.05}{n}\right)^{n\cdot2}$ dollars, after $t$ years $\left(1+\frac{0.05}{n}\right)^{nt}$ dollars.

Going one step further, if we let $A(t)$ denote the value of the investment
of one dollar earning $5\%$ annually, compounded $n$ times per year,
after $t$
years,  we have
$$
A(t)=\left(1+\frac{0.05}{n}\right)^{nt},\ t=0,1,2,\ldots
$$

What if the investment was compounded continuously?  What does this
mean, and what would the value $A(t)$ be?  To examine this, it will be
convenient to write $A(t)$ as
\begin{align*}
  A(t)&=\left(1+\frac{0.05}{n}\right)^{nt}\\
      &=\left(\left(1+\frac{0.05}{n}\right)^{\frac{n}{0.05}}\right)^{0.05t}\\
      &=\left(\left(1+\frac{1}{\frac{n}{0.05}}\right)^{\frac{n}{0.05}}\right)^{0.05t}
\end{align*}
Let’s rewrite $m=\frac{n}{0.05}$, and consider what would happen if we compounded a very
large number of times (let $n$ be very large).  Table 4 looks at some
values of $\left(1+\frac{1}{m}\right)^m$, where $m$ is large.
\begin{table}[h]
  \centering
  \begin{center}
    \begin{tabular}[h]{|c|c|}\hline
      m & $\left(1+\frac{1}{m}\right)^m$\bigstrut{}\\\hline
     100  & 2.70481382942\bigstrut{}\\
     1000  & 2.71692393224\bigstrut{}\\
     10000  & 2.71814592683\bigstrut{}\\
     100000  & 2.71826823717\bigstrut{}\\
     1000000  & 2.71828046932\bigstrut{}\\
     10000000  & 2.71828169255\bigstrut{}\\\hline
    \end{tabular}
  \end{center}
\end{table}
Does this look familiar to you.  It would appear from Table 4, that
$\left(1+\frac{1}{m}\right)^m$ approaches the number $e$.  We will show that this is the case later\notefrombob{This would make a
  nice problem in l”Hopital’s Rule section to re-examine. }.

All of this suggests that if the investment is being continually
compounded, then $A(t)=e^{0.05t}$. This fits in with the notion that
the investment grows continuously at a (nominal) rate of $5$\%. $A(t)$
must satisfy the IVP:
$$
\dfdx{A}{t}=0.05A,\ A(0)=1.
$$

\begin{embeddedproblem-enumerate}
  \label{problem:ContCompound1}
  \begin{enumerate}[label={\bf{}(\alph*)}]
  \item What would the effective yield be for a bond nominally rated
    at 5\% annually compounded continuously?  How does this compare to
    the effective yield of an investment compounded daily?
  \item Suppose we had two investments growing continually with
    nominal rates of 5\% and 10\% annually?  Would the effective yield
    of the second investment be twice that of the first?  Justify your
    answer.
  \end{enumerate}
\end{embeddedproblem-enumerate}

\begin{embeddedproblem}
  \label{PIC:DoubleInvestment}
  Of course, a natural question to ask is, how long will it take for
  my money to double?  Would an investment growing at 10\% double in
  half the time?

  Before we go any further, take a guess. Would an investment
  compounding continuously at a nominal rate of 10\% double in half
  the time as that of one growing at $5$\%.
\end{embeddedproblem}
To answer this question we will need to solve
$$
e^{0.05t}=2,
$$
for $t$ and compare this to the solution of
$$
e^{0.1t}=2.
$$

Loosely speaking, we need to find a way to ``undo'' the natural
exponential function.  This leads us to a discussion of the natural
logarithm function in the next section.
\end{myexample}



%%% Local Variables: 
%%% mode: latex
%%% outline-minor-mode: t
%%% eval:(flyspell-mode)
%%% TeX-master: "DiffCalc"
%%% End: 
